
\vspace{2cm}
\begin{center}
{\fontsize{14}{16.8}\selectfont\textbf{Bab VI \quad Kesimpulan dan Saran}}\\[1em]
\end{center}
\label{sec:kesimpulan-dan-saran}
\addcontentsline{toc}{section}{BAB VI KESIMPULAN DAN SARAN}
\setcounter{section}{6}
\setcounter{subsection}{0}
\setcounter{table}{0}
\setcounter{figure}{0}
\vspace{2em}

\subsection{Kesimpulan}
\label{subsec:kesimpulan}
\vspace{0.8em}

Penelitian ini berhasil mengembangkan kerangka kerja restorasi dokumen berorientasi HTR berbasis GAN yang mengatasi kesenjangan antara optimasi visual dan keterbacaan teks. Berdasarkan evaluasi empiris dan pengujian statistik, penelitian ini menyimpulkan dua kontribusi utama: (1) validasi strategi \textit{frozen recognizer} yang memberikan gradien sadar-teks stabil tanpa risiko keruntuhan modus (\textit{mode collapse}) yang umum terjadi pada pelatihan bersama ($p<0.001$), dan (2) identifikasi konfigurasi fungsi \textit{loss} multi-komponen yang mencapai keseimbangan Pareto optimal antara kualitas visual dan akurasi teks.

\vspace{0.8em}

Evaluasi pada 712 citra uji semi-sintetis dokumen paleografi ANRI abad ke-16--18 menunjukkan pengurangan \textit{Character Error Rate} (CER) yang signifikan dari 83.4\% menjadi 34.9\% (pengurangan relatif 58.2\%, $p<0.001$, \textit{Cohen's d}=0.85). Secara visual, model mempertahankan kualitas tinggi dengan PSNR 30.74 dB dan SSIM 0.987. Hasil CER pasca-restorasi (34.9\%) yang mendekati batas teoretis citra bersih (34.1\%) mengindikasikan bahwa hambatan utama kinerja HTR saat ini bukan lagi degradasi citra, melainkan keterbatasan intrinsik model pengenal dalam menangani ligatur paleografi kompleks. Validasi kualitatif pada 15 naskah ANRI autentik mengonfirmasi kemampuan generalisasi model pada degradasi historis nyata.

\vspace{0.8em}

Temuan kritis dari analisis komponen arsitektur adalah sebagai berikut:
\begin{enumerate}[leftmargin=1.5em, labelsep=0.5em, itemindent=0pt]
    \item Redundansi Diskriminator Dual-Modal: Arsitektur diskriminator dual-modal memberikan kontribusi marginal yang tidak signifikan ($\Delta$PSNR +0.28 dB, $p>0.05$). Hal ini membuktikan bahwa peran ``pengawasan tekstual'' telah diambil alih secara lebih efektif oleh mekanisme \textit{frozen recognizer} (melalui \textit{CTC loss}). Berdasarkan prinsip parsimoni, arsitektur diskriminator CNN tunggal (visual) direkomendasikan sebagai standar efisiensi untuk implementasi masa depan.
    \item Dominasi Protokol Pelatihan: Keberhasilan restorasi lebih ditentukan oleh strategi pelatihan (penyeimbangan dan pembekuan bobot) daripada kompleksitas arsitektur jaringan. \textit{Curriculum learning} terbukti bermanfaat untuk stabilitas pada fase awal (\textit{warmup}), meskipun dampaknya terhadap metrik performa akhir dapat diabaikan.
    \item Inefisiensi Komponen Fitur Rekognisi: Komponen \textit{recognition feature loss} (RecFeat) terbukti tidak memberikan peningkatan statistik pada akurasi teks dalam studi ablasi. Penggunaannya dalam model produksi saat ini diidentifikasi sebagai artefak eksperimental yang dapat dieliminasi untuk efisiensi komputasi tanpa mengorbankan performa.
    \item Validasi Multimetrik: Korelasi lemah antara PSNR dan CER ($r \approx -0.15$) dibandingkan dengan korelasi kuat antara SNR dan peningkatan CER ($r = -0.963$) menegaskan bahwa metrik visual semata tidak cukup untuk mengevaluasi efektivitas restorasi dokumen sejarah.
\end{enumerate}

\vspace{0.8em}

\subsection{Keterbatasan}
\label{subsec:keterbatasan}
\vspace{0.8em}

Keterbatasan penelitian yang teridentifikasi meliputi:
\begin{enumerate}[leftmargin=1.5em, labelsep=0.5em, itemindent=0pt]
    \item Degradasi Ekstrem: Model menghadapi tantangan pada dokumen dengan kontras sangat rendah ($<30$) dan pemudaran berat (SNR $< 10$ dB), di mana pendekatan konservatif model cenderung menghasilkan peningkatan minimal untuk menghindari halusinasi fitur.
    \item Keterbatasan Pengenal HTR: Analisis kegagalan menunjukkan bahwa 62.9\% sisa kesalahan (\textit{residual error}) disebabkan oleh ligatur paleografi kompleks yang gagal dikenali oleh model HTR, bukan karena kegagalan restorasi visual.
    \item Bias Aksara: Penelitian dibatasi pada aksara Latin paleografi (Bahasa Belanda Kuno), sehingga generalisasi ke aksara lain (Arab, Melayu, Jawa Kuno) belum tervalidasi.
    \item Kesenjangan Domain: Meskipun validasi kualitatif menunjukkan hasil positif, ketiadaan data \textit{ground truth} berpasangan untuk dokumen historis nyata membatasi evaluasi kuantitatif yang presisi pada data riil.
    \item Kompleksitas Komputasi: Durasi pelatihan (48 jam GPU untuk 50 \textit{epoch}) masih relatif tinggi untuk penerapan praktis yang membutuhkan pelatihan ulang berkala pada koleksi baru.
\end{enumerate}

\vspace{0.8em}

\subsection{Saran}
\label{subsec:saran}
\vspace{0.8em}

Berdasarkan temuan dan keterbatasan penelitian, saran diajukan bagi dua aspek: pengembangan keilmuan dan penerapan praktis di lapangan.

\subsubsection{Saran untuk Penelitian Lanjutan}
\label{subsubsec:saran-penelitian}
\vspace{0.8em}
Untuk mengatasi keterbatasan teknis dan memperluas cakupan keilmuan, penelitian mendatang disarankan berfokus pada:

\begin{enumerate}[leftmargin=1.5em, labelsep=0.5em, itemindent=0pt]
\item \textit{Fine-Tuning} Pengenal HTR pada Data Paleografi: Mengingat hambatan utama adalah pengenal (\textit{recognizer}), prioritas riset selanjutnya adalah melakukan \textit{fine-tuning} atau adaptasi domain pada model HTR menggunakan korpus spesifik paleografi ligatur kompleks untuk menurunkan batas bawah kesalahan teoretis.

\item Arsitektur Diskriminator yang Lebih Efisien: Mengganti desain dual-modal dengan diskriminator visual tunggal yang lebih ringan (\textit{lightweight}), serta mengeksplorasi fungsi \textit{loss} alternatif untuk menggantikan komponen RecFeat yang terbukti redundan.

\item Pemodelan Degradasi Berbasis Fisika: Pengembangan model degradasi kimia historis autentik (korosi tinta \textit{iron gall}, \textit{foxing} biologis) dan mekanisme augmentasi data yang lebih agresif untuk meningkatkan ketangguhan model terhadap kasus degradasi ekstrem.

\item Estimasi Ketidakpastian (\textit{Uncertainty Estimation}): Mengintegrasikan metode Bayesian atau \textit{ensemble} untuk memberikan skor kepercayaan pada setiap wilayah yang direstorasi, sehingga arsiparis dapat memverifikasi secara selektif bagian yang memiliki ketidakpastian tinggi.

\item Pembangunan Tolok Ukur (\textit{Benchmark}) Nasional: Pembangunan dataset standar dokumen paleografi Indonesia dengan anotasi \textit{ground truth} berkualitas tinggi (melalui kolaborasi ahli paleografi) untuk memungkinkan evaluasi kuantitatif yang objektif pada data riil.
\end{enumerate} 

\vspace{0.8em}

\subsubsection{Saran untuk Implementasi Praktis}
\label{subsubsec:saran-praktis}
\vspace{0.8em}
Bagi lembaga kearsipan (seperti ANRI) yang hendak mengadopsi teknologi ini, disarankan:

\begin{enumerate}[leftmargin=1.5em, labelsep=0.5em, itemindent=0pt]
\item Adopsi Bertahap: Implementasi sistem restorasi dimulai pada koleksi prioritas dengan tingkat degradasi menengah hingga berat, di mana model terbukti memberikan peningkatan keterbacaan paling signifikan.
\item Verifikasi Hibrida: Menggunakan keluaran model sebagai draf awal transkripsi untuk mempercepat kerja paleografer, namun tetap mempertahankan verifikasi manusia untuk dokumen bernilai tinggi guna menjamin autentisitas.
\item Standarisasi Digitalisasi: Penerapan standar pengambilan citra yang konsisten (resolusi, pencahayaan) untuk meminimalkan variabilitas domain yang dapat menurunkan kinerja model otomatis.
\end{enumerate}

\vspace{0.8em}

\subsection{Penutup}
\label{subsec:penutup}
\vspace{0.8em}

Penelitian ini menegaskan bahwa dalam pengembangan sistem pemulihan dokumen historis, sinergi antara pemrosesan citra dan pengenalan pola teks adalah kunci utama. Temuan bahwa protokol pelatihan yang tepat (\textit{frozen strategy}) lebih krusial daripada kompleksitas arsitektur memberikan panduan efisiensi bagi pengembangan sistem masa depan. Dengan kemampuan memulihkan keterbacaan hingga mendekati kondisi dokumen bersih, kerangka kerja ini menawarkan kontribusi nyata bagi upaya pelestarian warisan budaya digital dan demokratisasi akses terhadap informasi sejarah bangsa.

